\section{Conclusion}

John 6 supports a robust Eucharistic reading because its own language
interweaves faith with flesh given, true food and drink, mutual abiding, and
resurrection; early liturgical, homiletic, and doctrinal witnesses receive it
that way. The same tradition refuses crude physicalism by locating the gift in
the incarnate and life-giving Word and by insisting on spiritual, faithful,
ecclesial reception.

The Nativity conjunction is also older than modern devotional shorthand.
Jerome's Paula and Ambrose explicitly join Bethlehem, ``house of bread,'' to
the Bread from heaven; Chrysostom joins the manger-body to the Eucharistic
Body. The larger picture---Bethlehem, feeding trough, cloth-wrapped child, and
communion---is best presented as a canonical and traditional
theological synthesis. Its parts have ancient witnesses; its most vivid
all-in-one patristic proof-text presently rests on disputed attribution.

That distinction does not impoverish the image. It keeps the image from asking
history to say more than history says.

